\begin{comment}
\section{Formulácia problému, ciele a motivácia projektu}

Accurate and timely information about hazards found in food is an important topic due how to it can directly affect public health. At the same time, this information is also very relevant for food producers as food recalls can cause substantial financial damage. Automated systems that scan relevant food incident reports and extract the hazard and product affected could assist experts and food authorities in monitoring the presence of dangerous substances within products. The SemEval-2025 Task 9 challenge \cite{randl-etal-2025-semeval} focuses directly on this topic, with its goal being to assess how well machine learning models can extract data about hazards found in food from noisy, real-world recall announcements published by official food authorities, mainly in the United States, the United Kingdom, Canada and Australia. 

The challenge is divided into two sub-tasks. In the first sub-task (ST1) the goal is to correctly identify the broader category of both the hazard and the relevant product with there being 10 and 22 distinct categories of hazard and product categories respectively \cite{randl-etal-2025-semeval}. In the second, harder sub-task (ST2) the goal is to correctly identify the precise hazard and product affected with there being provided 128 labels for hazards and 1,142 labels for products \cite{randl-etal-2025-semeval}. The second sub-task was introduced as a way to broadly elaborate on why a report was classified in a certain category e.g. ``This report concerns allergens (hazard category) because it mentions pecan nut (hazard)''. Despite this fact, 21 out of the 27 submitted solutions preferred to create separate pipelines for each of the two sub-tasks, with only 6 solutions choosing to exploit the correlation between the two sub-tasks \cite{randl-etal-2025-semeval}.

As part of this challenge three baseline solutions were provided : a TF-IDF-based solution, a BERT-based solution and a CICLe-based \cite{randl-etal-2024-cicle} solution \cite{randl-etal-2025-semeval}. Our team decided to focus on the BERT based solution, with the goal of stage 1 of this project being to reproduce the metrics found in the official document. This solution will be used on both sub-tasks using the official data splits for training, validation and testing.

\end{comment}

\section{Formulácia problému, ciele a motivácia projektu}

Presné a včasné informácie o nebezpečných látkach v potravinách sú dôležitou témou, pretože môžu priamo ovplyvniť verejné zdravie. Zároveň sú tieto informácie veľmi dôležité aj pre výrobcov potravín, keďže sťahovanie potravín z trhu môže spôsobiť značné finančné škody. Automatizované systémy, ktoré prehľadávajú relevantné správy o potravinových incidentoch a extrahujú z nich nebezpečnú látku a dotknutý produkt, by mohli pomôcť expertom a potravinovým úradom široko monitorovať výskyt nebezpečných látok v produktoch. Súťaž SemEval-2025 Task 9 \cite{randl-etal-2025-semeval} sa priamo zameriava na túto tému. Jej cieľom je posúdiť, ako dobre dokážu modely strojového učenia extrahovať údaje o nebezpečných látkach v potravinách zo zašumených reálnych oznámení o sťahovaní produktov z trhu, ktoré zverejňujú oficiálne potravinové úrady, najmä v Spojených štátoch, Spojenom kráľovstve, Kanade a Austrálii.

Súťaž je rozdelená na dve podúlohy. V prvej podúlohe (ST1) je cieľom správne určiť širšiu kategóriu nebezpečnej látky (hazardu) aj príslušného produktu, pričom existuje 10 kategórií hazardov a 22 kategórií produktov \cite{wan-etal-2025-pateam}. V druhej, náročnejšej podúlohe (ST2) je cieľom správne určiť presný hazard a dotknutý produkt, pričom je k dispozícii 128 označení (labelov) pre hazardy a 1\,142 pre produkty \cite{wan-etal-2025-pateam}. Druhá podúloha bola zavedená ako spôsob, ako v širšom zmysle zdôvodniť, prečo bola správa zaradená do určitej kategórie, napr. „Táto správa sa týka allergens (kategória hazardu), pretože spomína pecan nut (hazard).“ Napriek tomu 21 z 27 odovzdaných riešení uprednostnilo samostatné pipeline pre každú z podúloh a iba 5 riešení sa rozhodlo využiť koreláciu medzi nimi \cite{randl-etal-2025-semeval}.

%V rámci súťaže boli poskytnuté tri referenčné (baseline) riešenia: riešenie založené na TF-IDF, riešenie založené na modeli BERT a riešenie založené na metóde CICLe \cite{randl-etal-2024-cicle, randl-etal-2025-semeval}. Náš tím sa rozhodol zamerať na riešenie založené na metóde CICLe, pričom cieľom 1.~etapy projektu je implementácia tohto riešenia a porovnanie jeho výkonu s zverejnenými metrikami ostatných (dvoch) baseline riešení. 

Cieľom tejto etapy projektu je vytvoriť vhodné baseline riešenie, ktoré bude zohľadňovať trendy nájdené v ostatných riešeniach odovzdaných do tejto súťaže. V rámci súťaže boli poskytnuté tri referenčné riešenia: riešenie založené na TF-IDF, riešenie založené na modeli BERT a riešenie založené na metóde CICLe \cite{randl-etal-2024-cicle, randl-etal-2025-semeval}. Naše baseline riešenie bude následne porovnané s týmito tromi referenčnými riešeniami na základe oficiálnej metriky a zároveň bude ukázané, kde by sa umiestnilo v oficiálnom rebríčku.      